% !TEX TS-program = xelatex
% !BIB program = bibtex
% !TeX spellcheck = ru_RU
% !TEX root = talk.tex

\documentclass
  [ russian
  , aspectratio=1610 % Для защит онлайн лучше использовать разрешение не 4х3
  ] {beamer}
\usepackage{blkarray}
\usepackage{listings}


% !TeX spellcheck = ru_RU
% !TEX root = vkr.tex
% Опциональные добавления используемых пакетов. Вполне может быть, что они вам не понадобятся, но в шаблоне приведены примеры их использования.
\usepackage{tikz} % Мощный пакет для создание рисунков, однако может очень сильно замедлять компиляцию
\usetikzlibrary{decorations.pathreplacing,calc,shapes,positioning,tikzmark}

% Библиотека для TikZ, которая генерирует отдельные файлы для каждого рисунка
% Позволяет ускорить компиляцию, однако имеет свои ограничения
% Например, ломает пример выделения кода в листинге из шаблона
% \usetikzlibrary{external}
% \tikzexternalize[prefix=figures/]

\newcounter{tmkcount}

\usepackage{fancyvrb} % для verbatim переноса строки

% --- переносы в коде/JSON ---
\usepackage{listings}
\lstset{
  basicstyle=\ttfamily\small,
  breaklines=true,           % перенос длинных строк
  breakatwhitespace=false,   % переносить можно в любом месте
  columns=fullflexible,      % корректная ширина моношрифта
  keepspaces=true,           % сохранять пробелы
  showstringspaces=false
}
% (необязательно) Определение простого языка JSON для подсветки
\lstdefinelanguage{json}{
  string=[s]{"}{"},
  morestring=[s]{'}{'}, % если вдруг одинарные кавычки
  comment=[l]{//},
}

\tikzset{
    use tikzmark/.style={
            remember picture,
            overlay,
            execute at end picture={
                    \stepcounter{tmkcount}
                },
        },
    tikzmark suffix={-\thetmkcount}
}

\usepackage{booktabs} % Пакет для верстки "более книжных" таблиц, вполне годится для оформления результатов
% В шаблоне есть команда \multirowcell, которой нужен этот пакет.
\usepackage{multirow}
\usepackage{siunitx} % для таблиц с единицами измерений

% Для названий стоит использовать \textsc{}
\newcommand{\OCaml}{\textsc{OCaml}}
\newcommand{\miniKanren}{\textsc{miniKanren}}
\newcommand{\BibTeX}{\textsc{BibTeX}}
\newcommand{\vsharp}{\textsc{V$\sharp$}}
\newcommand{\fsharp}{\textsc{F$\sharp$}}
\newcommand{\csharp}{\textsc{C\#}}
\newcommand{\GitHub}{\textsc{GitHub}}
\newcommand{\SMT}{\textsc{SMT}}

\definecolor{eclipseGreen}{RGB}{63,127,95}
% \lstdefinelanguage{ocaml}{
% keywords={@type, function, fun, let, in, match, with, when, class, type,
% nonrec, object, method, of, rec, repeat, until, while, not, do, done, as, val, inherit, and,
% new, module, sig, deriving, datatype, struct, if, then, else, open, private, virtual, include, success, failure,
% lazy, assert, true, false, end},
% sensitive=true,
% commentstyle=\small\itshape\ttfamily,
% keywordstyle=\ttfamily\bfseries, %\underbar,
% identifierstyle=\ttfamily,
% basewidth={0.5em,0.5em},
% columns=fixed,
% fontadjust=true,
% literate={->}{{$\to$}}3 {===}{{$\equiv$}}1 {=/=}{{$\not\equiv$}}1 {|>}{{$\triangleright$}}3 {\\/}{{$\vee$}}2 {/\\}{{$\wedge$}}2 {>=}{{$\ge$}}1 {<=}{{$\le$}} 1,
% morecomment=[s]{(*}{*)}
% }

\makeatletter
\@ifclassloaded{beamer}{
    %%% Обязательные пакеты
    %% Beamer
    \usepackage{beamerthemesplit}
    \usetheme{SPbGU}
    \beamertemplatenavigationsymbolsempty
    \usepackage{appendixnumberbeamer}

    %% Локализация
    \usepackage{fontspec}
    \setmainfont{CMU Serif}
    \setsansfont{CMU Sans Serif}
    \setmonofont{CMU Typewriter Text}
    %\setmonofont{Fira Code}[Contextuals=Alternate,Scale=0.9]
    %\setmonofont{Inconsolata}
    \usepackage{polyglossia}
    \setmainlanguage{russian}
    \setotherlanguage{english}

    %% Графика
    \usepackage{pdfpages} % Позволяет вставлять многостраничные pdf документы в текст

    % Математические окружения с русским названием
    \newtheorem{rutheorem}{Теорема}
    \newtheorem{ruproof}{Доказательство}
    \newtheorem{rudefinition}{Определение}
    \newtheorem{rulemma}{Лемма}
    \usepackage{fancyvrb}
}
{}
\makeatother

\usepackage[autostyle]{csquotes} % Правильные кавычки в зависимости от языка
\usepackage{totcount}
\usepackage{setspace}
\usepackage{amsmath, amsfonts, amssymb, amsthm, mathtools} % "Адекватная" работа с математикой в LaTeX


\makeatletter

%!TEX root = vkr.tex

%% Параметры заполнения титульного листа
\usepackage{xkeyval}

%% Русскоязычный вариант
\define@key[ru]{mytitle}{chair}{\def\my@title@chair@ru{#1}}
\define@key[ru]{mytitle}{title}{\def\my@title@title@ru{#1}}
\define@key[ru]{mytitle}{group}{\def\my@title@group@ru{#1}}
\define@key[ru]{mytitle}{author}{\def\my@title@author@ru{#1}}
\define@key[ru]{mytitle}{supervisor}{\def\my@title@supervisor@ru{#1}}
\define@key[ru]{mytitle}{supervisorPosition}{\def\my@title@supervisorPosition@ru{#1}}
\define@key[ru]{mytitle}{reviewer}{\def\my@title@reviewer@ru{#1}}
\define@key[ru]{mytitle}{reviewerPosition}{\def\my@title@reviewerPosition@ru{#1}}
\define@key[ru]{mytitle}{consultant}{\def\my@title@consultant@ru{#1}}
\define@key[ru]{mytitle}{consultantPosition}{\def\my@title@consultantPosition@ru{#1}}
\define@key[ru]{mytitle}{year}{\def\my@title@year@ru{#1}}
\define@key[ru]{mytitle}{specialty}{\def\my@title@specialty@ru{#1}}
\define@key[ru]{mytitle}{programme}{\def\my@title@programme@ru{#1}}
\define@key[ru]{mytitle}{profile}{\def\my@title@profile@ru{#1}}
\define@choicekey*[ru]{mytitle}{type}{coursework,practice,prediploma,master,bachelor,production}{\def\my@title@type@ru{#1}}
\define@choicekey*[ru]{mytitle}{kind}{solution,experiment,production,comparison,theoretical}{\def\my@title@kind@ru{#1}}

%% Англоязычный вариант
\define@key[en]{mytitle}{chair}{\def\my@title@chair@en{#1}}
\define@key[en]{mytitle}{title}{\def\my@title@title@en{#1}}
\define@key[en]{mytitle}{group}{\def\my@title@group@en{#1}}
\define@key[en]{mytitle}{author}{\def\my@title@author@en{#1}}
\define@key[en]{mytitle}{supervisor}{\def\my@title@supervisor@en{#1}}
\define@key[en]{mytitle}{supervisorPosition}{\def\my@title@supervisorPosition@en{#1}}
\define@key[en]{mytitle}{reviewer}{\def\my@title@reviewer@en{#1}}
\define@key[en]{mytitle}{reviewerPosition}{\def\my@title@reviewerPosition@en{#1}}
\define@key[en]{mytitle}{consultant}{\def\my@title@consultant@en{#1}}
\define@key[en]{mytitle}{consultantPosition}{\def\my@title@consultantPosition@en{#1}}
\define@key[en]{mytitle}{year}{\def\my@title@year@en{#1}}
\define@key[en]{mytitle}{specialty}{\def\my@title@specialty@en{#1}}
\define@key[en]{mytitle}{programme}{\def\my@title@programme@en{#1}}
\define@key[en]{mytitle}{profile}{\def\my@title@profile@en{#1}}
\define@choicekey*[en]{mytitle}{type}{coursework,practice,prediploma,master,bachelor}{\def\my@title@type@en{#1}}
\define@choicekey*[en]{mytitle}{kind}{solution,experiment,production,comparison,theoretical}{\def\my@title@kind@en{#1}}

\newcommand{\filltitle}[2]{
    %% Значения по умолчанию для обоих языков
    \ifthenelse{\equal{#1}{ru}}
    {
        \presetkeys[#1]{mytitle}{
            year = {\the\year},
            type = {practice},
            reviewer = {},
            consultant = {},
            profile = {}
        }{}
    }
    {}
    \ifthenelse{\equal{#1}{en}}
    {
        \presetkeys[#1]{mytitle}{
            year = {\the\year},
            type = {practice},
            reviewer = {},
            consultant = {},
            profile = {}
        }{}
    }
    {}
    \setkeys[#1]{mytitle}{#2}
}


% !TeX spellcheck = ru_RU
% !TEX root = vkr.tex

%% Если что-то забыли, при компиляции будут ошибки Undefined control sequence \my@title@<что забыли>@ru
%% Если англоязычная титульная страница не нужна, то ее можно просто удалить.
\filltitle{ru}{
    %% Актуально только для курсовых/практик. ВКР защищаются не на кафедре а в ГЭК по направлению,
    %%   и к моменту защиты вы будете уже не в группе.
    chair              = {Кафедра информатики},
    group              = {23.Б16-мм},
    %
    %% Макрос filltitle ненавидит пустые строки, поэтому обязателен хотя бы символ комментария на строке
    %% Актуально всем.
    title              = {Реализация и внедрение архитектуры Gated Recurrent Unit в пакет DEGANN},
    %
    %% Здесь указывается тип работы. Возможные значения:
    %%   production - производственная практика;
    %%   coursework - отчёт по курсовой работе (ОБРАТИТЕ ВНИМАНИЕ, у техпрога и ПИ нет курсовых, только практики);
    %%   practice - отчёт по учебной практике;
    %%   prediploma - отчёт по преддипломной практике;
    %%   master - ВКР магистра;
    %%   bachelor - ВКР бакалавра.
    type               = {practice},
    %
    %% Здесь указывается вид работы. От вида работы зависят критерии оценивания.
    %%   solution - «Решение». Обучающемуся поручили найти способ решения проблемы в области разработки программного обеспечения или теоретической информатики с учётом набора ограничений.
    %%   experiment - «Эксперимент». Обучающемуся поручили изучить возможности, достоинства и недостатки новой технологии, платформы, языка и т. д. на примере какой-то задачи.
    %%   production - «Производственное задание». Автору поручили реализовать потенциально полезное программное обеспечение.
    %%   comparison - «Сравнение». Обучающемуся поручили сравнить несколько существующих продуктов и/или подходов.
    %%   theoretical - «Теоретическое исследование». Автору поручили доказать какое-то утверждение, исследовать свойства алгоритма и т.п., при этом не требуя написания кода.
    kind               = {production},
    %
    author             = {Мурадян Денис Степанович},
    %
    %% Актуально только для ВКР. Указывается код и название направления подготовки. Типичные примеры:
    %%   02.03.03 \enquote{Математическое обеспечение и администрирование информационных систем}
    %%   02.04.03 \enquote{Математическое обеспечение и администрирование информационных систем}
    %%   09.03.04 \enquote{Программная инженерия}
    %%   09.04.04 \enquote{Программная инженерия}
    %% Те, что с 03 в середине --- бакалавриат, с 04 --- магистратура.
    specialty          = {02.03.03 \enquote{Математическое обеспечение и администрирование информационных систем}},
    %
    %% Актуально только для ВКР. Указывается шифр и название образовательной программы. Типичные примеры:
    %%   СВ.5162.2020 \enquote{Технологии программирования}
    %%   СВ.5080.2020 \enquote{Программная инженерия}
    %%   ВМ.5665.2022 \enquote{Математическое обеспечение и администрирование информационных систем}
    %%   ВМ.5666.2022 \enquote{Программная инженерия}
    %% Шифр и название программы можно посмотреть в учебном плане, по которому вы учитесь.
    %% СВ.* --- бакалавриат, ВМ.* --- магистратура. В конце --- год поступления (не обязательно ваш, если вы были в академе/вылетали).
    programme          = {СВ.5162.2020 \enquote{Технологии программирования}},
    %
    %% Актуально всем.
    %% Должно умещаться в одну строчку, допускается использование сокращений, но без переусердствования,
    %% короткая строка с большим количеством сокращений выглядит странно
    %supervisorPosition = {проф. кафeдры системного программирования, д.ф.-м.н.,}, % Терехов А. Н.
    %supervisorPosition = {ст. преподаватель кафедры ИАС, к.~ф.-м.~н. (если есть),}, % Смирнов К. К.
    supervisorPosition = {доцент кафедры системного программирования, к.ф.-м.н.,},
    supervisor         = {Гориховский~В.~И.},
    %
    %% Актуально только для практик и курсовых. Если консультанта нет или он совпадает с научником, закомментировать или удалить вовсе.
    consultantPosition = {Преподаватель СПбГУ,},
    consultant         = {Алимов~П.~Г.},
}

% Английский титульник нужен только для ВКР, остальные виды работ могут его смело игнорировать.
\filltitle{en}{
    chair              = {Advisor's chair},
    group              = {ХХ.BХХ-mm},
    title              = {Template for SPbU qualification works},
    type               = {bachelor},
    author             = {FirstName Surname},
    %
    %% Possible choices:
    %%   02.03.03 \foreignquote{english}{Software and Administration of Information Systems}
    %%   02.04.03 \foreignquote{english}{Software and Administration of Information Systems}
    %%   09.03.04 \foreignquote{english}{Software Engineering}
    %%   09.04.04 \foreignquote{english}{Software Engineering}
    %% Те, что с 03 в середине --- бакалавриат, с 04 --- магистратура.
    specialty          = {02.03.03 \foreignquote{english}{Software and Administration of Information Systems}},
    %
    %% Possible choices:
    %%   СВ.5162.2020 \foreignquote{english}{Programming Technologies}
    %%   СВ.5080.2020 \foreignquote{english}{Software Engineering}
    %%   ВМ.5665.2022 \foreignquote{english}{Software and Administration of Information Systems}
    %%   ВМ.5666.2022 \foreignquote{english}{Software Engineering}
    programme          = {СВ.5162.2020 \foreignquote{english}{Programming Technologies}},
    %
    %% Note that common title translations are:
    %%   кандидат наук --- C.Sc. (NOT Ph.D.)
    %%   доктор ... наук --- Sc.D.
    %%   доцент --- docent (NOT assistant/associate prof.)
    %%   профессор --- prof.
    supervisorPosition = {Sc.D, prof.},
    supervisor         = {S.S. Supervisor},
    %
    consultantPosition = {position at \foreignquote{english}{Company}, degree if present},
    consultant         = {C.C. Consultant},
    %
    reviewerPosition   = {position at \foreignquote{english}{Company}, degree if present},
    reviewer           = {R.R. Reviewer},
}


\newcommand{\academicGroup}{\my@title@group@ru}
\newcommand{\advisorChair}{\my@title@chair@ru}
% То, что в квадратных скобках, отображается внизу по центру каждого слайда.
\title[Короткое название]{\my@title@title@ru}
% То, что в квадратных скобках, отображается в левом нижнем углу.
\author[\my@title@author@ru]{\my@title@author@ru, группа \academicGroup}
\institute[СПбГУ]{}
\date[9 сентября 2024 г.]{}
\newcommand{\supervisor}{\my@title@supervisor@ru}
\newcommand{\supervisorPosition}{\my@title@supervisorPosition@ru}
\newcommand{\consultant}{\my@title@consultant@ru}
\newcommand{\consultantPosition}{\my@title@consultantPosition@ru}
\newcommand{\reviewer}{\my@title@reviewer@ru}
\newcommand{\reviewerPosition}{\my@title@reviewerPosition@ru}
\newcommand{\defenseYear}{\my@title@year@ru}

\makeatother
\begin{document}
{
\setbeamertemplate{footline}{}
% Лого университета или организации, отображается в шапке титульного листа
\begin{frame}
    \includegraphics[width=1.4cm]{figures/SPbGU_Logo.png}
    \vspace{-35pt}
    \hspace{-10pt}
    \begin{center}
        \begin{tabular}{c}
            \scriptsize{Санкт-Петербургский государственный университет} \\
            \scriptsize{\advisorChair}
        \end{tabular}
        \titlepage
    \end{center}

    \btVFill

    {\scriptsize
        % У научного руководителя должна быть указана научная степень
        \textbf{Научный руководитель:}  \supervisorPosition~\supervisor \\
        % Консультанта может и не быть. Должна быть указана должность или ученая степень
        \textbf{Консультант:}  \consultantPosition~\consultant \\
    }
    \makeatother
    \begin{center}
        \vspace{5pt}
        \scriptsize{Санкт-Петербург\\ \defenseYear}
    \end{center}
\end{frame}
}


\begin{frame}{Введение}
    \begin{itemize}
        \item Веб-парсинг: извлечение структурированных данных из веб-страниц
        \item Проблема: сайты бывают статическими и динамическими, для каждого нужен свой парсер и его постоянная поддержка
        \item Идея: автоматизация генерации parser-скриптов с помощью больших языковых моделей (LLM)
    \end{itemize}
\end{frame}

\begin{frame}{Существующие решения и их недостатки}
    \begin{itemize}
        \item \textbf{Статический парсинг (requests + BeautifulSoup)}
            \begin{itemize}
                \item + Лёгкость и скорость запуска
                \item – Не работает с динамическим JS-контентом, требует ручных селекторов
            \end{itemize}
        \item \textbf{Динамический парсинг (Selenium)}
            \begin{itemize}
                \item + Обходит любой JS
                \item – Тоже требует писать и поддерживать скрипты-парсеры под каждую структуру
            \end{itemize}
        \item \textbf{Общий недостаток}: отсутствует единый самонастраивающийся механизм — при изменении верстки ручной код ломается
    \end{itemize}
\end{frame}

\begin{frame}{LLM-подход и семантический кэш}
    \begin{itemize}
        \item \textbf{Structuring}: извлечение ответа на запрос пользователя напрямую из очищенного текста страницы
        \item \textbf{Codegen (двойное применение LLM)}:
            \begin{itemize}
                \item \emph{HintGen}: генерирует подсказки, селекторы и few-shot примеры
                \item \emph{CodeGen}: на основе подсказок формирует Python-скрипт-парсер и выполняет его
            \end{itemize}
        \item \textbf{Проблемы}: дороговизна запросов в LLM и время инференса ответа
        \item \textbf{Решение}: семантический кэш (SQLite + ChromaDB), для одинаковых ссылок и схожих embedding запросов - запросы в LLM не отправляются, а используются ранее кешированные парсеры
    \end{itemize}
\end{frame}

\begin{frame}{Постановка задачи}
    \textbf{Цель проекта:} оптимизировать сбор данных из открытых веб-источников за счёт автоматизации написания парсеров веб-страниц с помощью больших языковых моделей (LLM)

    \vspace{1em}
    \textbf{Задачи:}
    \begin{itemize}
        \item Определять тип страницы (статическая или динамическая)
        \item Реализовать две стратегии очистки HTML (Full / Lite)
        \item Внедрить режимы LLM-парсинга: Structuring и Codegen
        \item Построить семантический кэш (SQLite + ChromaDB)
        \item Разработать интерфейсы: REST API, веб-фронтенд, Gradio
        \item Провести экспериментальное исследование
    \end{itemize}
\end{frame}


\begin{frame}{Архитектура решения}
    \begin{center}
        \resizebox{\textwidth}{!}{%
        \begin{tikzpicture}[
            >=latex,
            node distance=0.8cm and 1.5cm,
            block/.style={
                draw=black,
                fill=blue!15,
                thick,
                rectangle,
                rounded corners,
                minimum height=2cm,
                minimum width=4cm,
                inner sep=5mm,
                align=center,
                font=\bfseries\Large
            },
            arrow/.style={
                ->,
                thick
            }
        ]
            % Основная линия
            \node[block]                       (input)   {Ввод URL\\и запроса};
            \node[block,right=of input]        (detect)  {Детектор\\динамичности};
            \node[block,right=of detect]       (fetch)   {Fetcher\\(requests / Selenium)};
            \node[block,right=of fetch]        (clean)   {Очистка\\(Full / Lite)};
            \node[block,right=of clean]        (dispatch){Диспетчер\\режимов};
            \node[block,right=of dispatch]     (cache)   {Семантический\\кэш};
            \node[block,right=of cache]        (output)  {Вывод JSON};

            % Ветви LLM
            \node[block,above=of dispatch]     (struct)  {Structuring\\(извлечение ответа)};
            \node[block,below=of dispatch]     (codegen) {Codegen\\(HintGen → CodeGen)};

            % Стрелки
            \draw[arrow] (input)    -- (detect)
                          (detect)  -- (fetch)
                          (fetch)   -- (clean)
                          (clean)   -- (dispatch)
                          (dispatch)-- (cache)
                          (cache)   -- (output);

            \draw[arrow] (dispatch) -- (struct) -- (cache);
            \draw[arrow] (dispatch) -- (codegen) -- (cache);
        \end{tikzpicture}
        }
    \end{center}
\end{frame}




\begin{frame}{Результаты экспериментального исследования}
    \begin{table}
        \centering
        \small
        \begin{tabular}{lcccc}
            \toprule
            \textbf{Сайт / Запрос} & \textbf{Codegen (cold)} & \textbf{Codegen (warm)} & \textbf{Structuring} \\
            \midrule
            \textbf{Gismетео} & & & & \\
            «Текущая температура в СПБ»      & 23.28 & 5.61 & 8.62 \\
            \addlinespace
            \textbf{СПбГУ} & & & & \\
            «Выведи последние новости СПбГУ»  & 40.91 & 6.09 & 7.34 \\
            \addlinespace
            \textbf{Яндекс.Финанс} & & & & \\
            «Курс доллара к рублю»           & 17.59 & 4.61 & 12.23 \\
            \bottomrule
        \end{tabular}
    \end{table}
\end{frame}


\begin{frame}[fragile]{Примеры JSON-ответов}
    \begin{block}{Gismетео: «Текущая температура в Санкт-Петербурге»}
\begin{lstlisting}[language=json,breaklines,keepspaces]
{
  "query_data": "Текущая температура в Санкт-Петербурге: +11°C\nТемпература по ощущению: +10°C\n"
}
\end{lstlisting}
    \end{block}

    \begin{block}{Яндекс.Финанс: «Курс доллара США к рублю»}
\begin{lstlisting}[language=json,breaklines,keepspaces]
{
  "query_data": "Курс доллара США к рублю: 78,62 руб\nИзменение курса: −2,88 руб (−3,53%)\nИсточник: ЦБ РФ\n"
}
\end{lstlisting}
    \end{block}
\end{frame}

\begin{frame}{Итоги работы}
    \begin{itemize}
        \item Разработана модульная система для парсинга статических и динамических веб-страниц
        \item Реализованы два режима LLM-генерации парсеров:
            \begin{itemize}
                \item прямое извлечение ответа из текста (Structuring)
                \item генерация Python-скрипта с помощью подсказок (Codegen)
            \end{itemize}
        \item Внедрён семантический кэш для повторного использования сгенерированных скриптов
        \item Предоставлены разные интерфейсы: библиотека, REST API, веб-приложение и Gradio
        \item Экспериментально подтверждена эффективность подхода на реальных сайтах
    \end{itemize}
\end{frame}


\begin{frame}
    \begin{center}
        \includegraphics[width=6cm]{figures/qr.png}
        \vspace{1em}

        \scriptsize{Воспользоваться проектом можно, отсканировав QR-код для перехода на Hugging Face Space}
    \end{center}
\end{frame}

\end{document}

